\section{Background and Related Work}
\label{sec:related}

This thesis sits at the intersection of two research strands, the automated
classification of astronomical spectra and the use of variational quantum circuits
for machine learning. This section reviews the work that the thesis builds on and
positions our contribution against it.

\subsection{Deep Learning for Spectral Classification}
\label{sec:rel_spectra}

The Sloan Digital Sky Survey has produced more than two million optical spectra and
established large-scale spectroscopy as a routine source of
data~\cite{york2000sdss}. The classification and redshift measurement of these
spectra were originally handled by template-fitting pipelines, in which an observed
spectrum is matched against a library of stellar, galaxy and quasar templates by a
least-squares fit. The SDSS-III pipeline of Bolton et al.~\cite{bolton2012boss} is a
mature example and remains a standard reference for the survey. Template fitting is
accurate on clean spectra but scales poorly to the volume of modern surveys and
degrades on noisy or atypical objects.

Deep learning has since become the dominant approach. Sharma et
al.~\cite{sharma2020cnn} showed that convolutional neural networks classify stellar
spectra directly from flux and that deeper networks exploit finer spectral detail
than shallow models. More recently the Galaxy Spectra Neural Network of Zhong et
al.~\cite{zhong2024gasnet}, GaSNet-II, performs joint classification and redshift
prediction across stars, galaxies and quasars and reaches roughly 99\% accuracy on
the coarse three-way problem while scaling to a customisable number of subclasses.
GaSNet-II is the natural reference point for the classical side of this work. We
benchmark our own classical baselines against the level of performance it reports.
What this body of work establishes is that the coarse classification problem is
essentially solved by classical deep learning, which is precisely why we move to the
harder fine-grained subclasses where classical models still struggle.

\subsection{Variational Quantum Circuits}
\label{sec:rel_vqc}

A variational quantum circuit is a parameterised quantum circuit whose gate angles
are optimised by classical gradient descent, a paradigm introduced for near-term
hardware as quantum circuit learning by Mitarai et al.~\cite{mitarai2018qcl} and
surveyed more broadly in the reviews of quantum machine learning by Biamonte et
al.~\cite{biamonte2017qml} and of variational quantum algorithms by Cerezo et
al.~\cite{cerezo2021vqa}. The appeal is that even a small circuit acts on a Hilbert space whose dimension grows
as $2^n$ with the number of qubits, which raises the question of how expressive such
models really are and how their inputs should be encoded.

Two results inform the design choices made in this thesis. P{\'e}rez-Salinas et
al.~\cite{perezsalinas2020reuploading} introduced data re-uploading, in which the
classical input is re-injected at every layer of the circuit and showed that even a
single qubit can act as a universal classifier when its data is uploaded repeatedly.
Schuld et al.~\cite{schuld2021encoding} made the underlying mechanism precise by
showing that a variational quantum model is effectively a Fourier series in its
input, with the accessible frequencies fixed by the data-encoding gates. A Fourier
view is not unique to quantum circuits, since classical models can also be analysed in
frequency space. What is specific here is that the encoding gates bound the accessible
frequencies directly, which in turn bounds the functions the model can represent. Against
these expressivity results stands a trainability obstacle. McClean et
al.~\cite{mcclean2018barren} showed that deep, wide, randomly initialised circuits
suffer barren plateaus, where gradients vanish exponentially in the number of qubits
and training stalls. This trade-off between width, depth and trainability is exactly
what our architecture sweep probes.

\subsection{Quantum Kernels and Quantum Support Vector Machines}
\label{sec:rel_qsvm}

A separate line of work treats the quantum circuit not as a trainable classifier but
as a fixed feature map. Schuld and Killoran~\cite{schuld2019feature} and in the same
year, Havl{\'\i}{\v{c}}ek et al.~\cite{havlicek2019} showed that a quantum circuit can
embed classical data into a high-dimensional feature Hilbert space whose inner
products define a kernel, which a classical support vector machine then uses for
classification. This reframes the search for a quantum advantage as a question about
the kernel rather than about a trained network. Liu et al.~\cite{liu2021speedup} gave
a rigorous example of such an advantage by constructing a classification task, based
on the discrete logarithm problem, on which a quantum kernel provably outperforms any
classical learner. These guarantees are important, but they depend on carefully
constructed problem structure. No comparable advantage is known for generic
real-world data such as the spectra studied here. We therefore treat the kernel route
as an open empirical question. In Experiment 3 we evaluate a quantum support vector
machine on our own data in a few-shot regime, the setting where quantum kernels are
most often expected to help.

\subsection{Quantum Convolution}
\label{sec:rel_quanv}

Convolutional structure has been carried into the quantum setting in two distinct
ways. Cong et al.~\cite{cong2019qcnn} proposed the quantum convolutional neural
network, a circuit with a convolution-and-pooling structure that acts on quantum
states and trains with a number of parameters that grows only logarithmically with
the system size. Henderson et al.~\cite{henderson2020quanvolutional} instead
proposed the quanvolutional layer, in which a small fixed or trainable quantum
circuit is slid across a classical input like a convolutional kernel, transforming
local patches into feature maps that a classical network then processes. The
quanvolutional approach is the one most relevant to this thesis, since it applies a
small circuit repeatedly to a large classical signal rather than once at a bottleneck
and our noise-robustness experiment builds directly on it.

Both lines are taken up directly in this work, and in both cases we adapt them from the
two-dimensional image setting in which they were introduced to the one-dimensional
spectra studied here. Experiment~4 follows Henderson et
al.~\cite{henderson2020quanvolutional} with a quanvolutional filter slid along the
wavelength axis, and asks a question their original work does not, namely whether the
bounded, unitary nature of the quantum feature map yields any robustness to input noise
once the classical filter it replaces is held to the same parameter count. Experiment~5
implements the QCNN of Cong et al.~\cite{cong2019qcnn} in full, with weight-shared
$SU(4)$ convolutions and coherent pooling, and uses it as a negative control. The
architecture was designed for quantum data such as phases of matter, so applying it to a
classically measured spectrum removes the premise it was built on; what remains testable
is whether its inductive bias and its guaranteed absence of barren
plateaus~\cite{pesah2021absence} still buy anything once the input has to be embedded
classically. To our knowledge neither question has been examined on astronomical spectra.

\subsection{Positioning of this Thesis}
\label{sec:rel_position}

The work above leaves room for a focused study. Deep learning is the established tool
for spectral classification and variational quantum circuits are an actively studied
alternative, but comparing the two fairly is harder than it looks. We are not the first
to make this point. Schuld and Killoran~\cite{schuld2022advantage} argue that chasing a
demonstration of quantum advantage is often the wrong framing for quantum machine
learning research, and Bowles et al.~\cite{bowles2024benchmarking} show empirically that
many reported advantages do not survive comparison against properly tuned classical
models of matched capacity, and that benchmarking practice in the field is frequently
too weak to support the claims built on it. Our methodological stance is theirs, applied
to a specific scientific domain. In particular, a common pattern in applied quantum
machine learning places a powerful classical network in front of a small circuit and
credits the result to the quantum component without an equal-capacity classical control.
What this thesis adds is a systematic instance of such a comparison on a realistic
scientific signal, benchmarking hybrid variational circuits against parameter-matched
classical models on hard SDSS subclass tasks. Both sides use the same data and the same surrounding
pipeline. The classical control is built with the same care as the quantum model. We
study five settings, a variational circuit on the hardest binary pair, a hybrid
quantum head on a four-class task, a quantum support vector machine in a few-shot
regime, a quanvolutional filter under noise and a fully coherent QCNN as a negative
control.
